% Options for packages loaded elsewhere
\PassOptionsToPackage{unicode}{hyperref}
\PassOptionsToPackage{hyphens}{url}
%
\documentclass[
  11pt,
  a4paper,
]{article}
\usepackage{amsmath,amssymb}
\usepackage{iftex}
\ifPDFTeX
  \usepackage[T1]{fontenc}
  \usepackage[utf8]{inputenc}
  \usepackage{textcomp} % provide euro and other symbols
\else % if luatex or xetex
  \usepackage{unicode-math} % this also loads fontspec
  \defaultfontfeatures{Scale=MatchLowercase}
  \defaultfontfeatures[\rmfamily]{Ligatures=TeX,Scale=1}
\fi
\usepackage{lmodern}
\ifPDFTeX\else
  % xetex/luatex font selection
  \setmainfont[]{Liberation Serif}
  \setsansfont[]{Lato}
  \setmonofont[]{DejaVu Sans Mono}
\fi
% Use upquote if available, for straight quotes in verbatim environments
\IfFileExists{upquote.sty}{\usepackage{upquote}}{}
\IfFileExists{microtype.sty}{% use microtype if available
  \usepackage[]{microtype}
  \UseMicrotypeSet[protrusion]{basicmath} % disable protrusion for tt fonts
}{}
\makeatletter
\@ifundefined{KOMAClassName}{% if non-KOMA class
  \IfFileExists{parskip.sty}{%
    \usepackage{parskip}
  }{% else
    \setlength{\parindent}{0pt}
    \setlength{\parskip}{6pt plus 2pt minus 1pt}}
}{% if KOMA class
  \KOMAoptions{parskip=half}}
\makeatother
\usepackage{xcolor}
\usepackage[margin=22mm]{geometry}
\usepackage{longtable,booktabs,array}
\usepackage{calc} % for calculating minipage widths
% Correct order of tables after \paragraph or \subparagraph
\usepackage{etoolbox}
\makeatletter
\patchcmd\longtable{\par}{\if@noskipsec\mbox{}\fi\par}{}{}
\makeatother
% Allow footnotes in longtable head/foot
\IfFileExists{footnotehyper.sty}{\usepackage{footnotehyper}}{\usepackage{footnote}}
\makesavenoteenv{longtable}
\setlength{\emergencystretch}{3em} % prevent overfull lines
\providecommand{\tightlist}{%
  \setlength{\itemsep}{0pt}\setlength{\parskip}{0pt}}
\setcounter{secnumdepth}{5}
\usepackage{xcolor}
\definecolor{Ink}{HTML}{163043}
\definecolor{Accent}{HTML}{137A83}
\definecolor{Muted}{HTML}{556773}
\definecolor{Pale}{HTML}{EDF4F5}
\usepackage{titlesec}
\usepackage{fancyhdr}
\usepackage{needspace}
\usepackage{xurl}
\usepackage{hyperref}
\usepackage{microtype}
\usepackage{float}
\usepackage{fvextra}
\usepackage{etoolbox}
\usepackage{tikz}
\usetikzlibrary{positioning,arrows.meta}
\hypersetup{colorlinks=true,linkcolor=Ink,urlcolor=Accent,citecolor=Accent,pdftitle={Rineya Consequence World Model: Research and Implementation Plan},pdfauthor={Prepared for Cheng Li}}
\setlength{\emergencystretch}{3em}
\setlength{\parindent}{0pt}
\setlength{\parskip}{5.5pt plus 1pt}
\setlength{\headheight}{17pt}
\setlength{\headsep}{16pt}
\pagestyle{fancy}
\fancyhf{}
\fancyhead[L]{\sffamily\small\color{Muted}RINEYA / CONSEQUENCE WORLD MODEL}
\fancyhead[R]{\sffamily\small\color{Muted}RESEARCH EDITION 1.0}
\fancyfoot[L]{\sffamily\footnotesize\color{Muted}5 October 2026\quad |\quad Proposed native work; scoped CPU checks}
\fancyfoot[R]{\sffamily\small\color{Ink}\thepage}
\renewcommand{\headrulewidth}{0.3pt}
\renewcommand{\footrulewidth}{0pt}
\titleformat{\section}{\sffamily\LARGE\bfseries\color{Ink}}{\color{Accent}\thesection}{0.75em}{}
\titleformat{\subsection}{\sffamily\large\bfseries\color{Ink}}{\thesubsection}{0.65em}{}
\titleformat{\subsubsection}{\sffamily\normalsize\bfseries\color{Accent}}{\thesubsubsection}{0.6em}{}
\titlespacing*{\section}{0pt}{20pt}{12pt}
\titlespacing*{\subsection}{0pt}{15pt}{6pt}
\pretocmd{\section}{\Needspace{12\baselineskip}}{}{}
\let\RCWMoriginaltoc\tableofcontents
\renewcommand{\tableofcontents}{\RCWMoriginaltoc\tableofcontentsend}
\newcommand{\tableofcontentsend}{\clearpage}
\setcounter{tocdepth}{1}
\setcounter{secnumdepth}{2}
\AtBeginEnvironment{longtable}{\small\setlength{\tabcolsep}{5pt}\renewcommand{\arraystretch}{1.2}}
\AtBeginEnvironment{quote}{\color{Muted}}
\DefineVerbatimEnvironment{Highlighting}{Verbatim}{breaklines,commandchars=\\\{\},fontsize=\small}
\ifLuaTeX
  \usepackage{selnolig}  % disable illegal ligatures
\fi
\usepackage{bookmark}
\IfFileExists{xurl.sty}{\usepackage{xurl}}{} % add URL line breaks if available
\urlstyle{same}
\hypersetup{
  hidelinks,
  pdfcreator={LaTeX via pandoc}}

\author{}
\date{}

\begin{document}

\begin{titlepage}
\thispagestyle{empty}
\vspace*{0.45cm}
{\sffamily\bfseries\color{Accent}\large RINEYA RESEARCH / RCWM-001\par}
\vspace{1.25cm}
{\sffamily\bfseries\color{Ink}\fontsize{37}{43}\selectfont Actions have\par consequences.\par}
\vspace{0.6cm}
{\sffamily\color{Ink}\fontsize{21}{28}\selectfont A consequence world model\par for compact language intelligence\par}
\vspace{0.9cm}
{\large Research and implementation plan\par}
\vspace{0.25cm}
{\color{Muted}Typed action representations. Learned consequence embeddings.\par Goal-conditioned utility. Long-horizon and parallel decision-making.\par}
\vspace{0.9cm}
\noindent\colorbox{Pale}{\parbox{0.92\textwidth}{\vspace{0.25cm}\hspace{0.15cm}\begin{minipage}{0.86\textwidth}\sffamily\color{Ink}\textbf{The central test}\par
Can a learned consequence representation improve real action selection or reduce the labels needed for new goals, compared with a matched scalar evaluator?
\end{minipage}\vspace{0.25cm}}}
\vfill
{\sffamily\bfseries\color{Ink}Prepared for Cheng Li\par}
\vspace{0.2cm}
{\sffamily\color{Muted}5 October 2026 / Research edition 1.0\par}
\vspace{0.5cm}
{\small\color{Muted}43 primary-source references. A gated experimental program.\par Editable specifications, schemas, configurations and tested CPU reference code.\par No newly trained or qualified native Rineya model is claimed.\par}
\end{titlepage}
\setcounter{page}{1}

{
\setcounter{tocdepth}{1}
\tableofcontents
}
\section{Executive decision}\label{executive-decision}

\textbf{Build a consequence-conditioned decision model for the exact
Rineya 0.6B-class checkpoint. Start with a small, independently testable
predictor attached to frozen Rineya representations; advance to typed
consequence and action slots inside Rineya only when the predictive
representation improves real decisions.}

The motivating idea is to give actions a learned meaning in terms of
what they cause. Instead of training a network only to output a scalar
approval score, train it to predict a reusable representation of
consequences. A separate goal readout determines which consequences are
valuable. Attention can relate a proposed action to distant context,
prerequisites, concurrent work and delayed outcomes. This is the user's
central proposal, formalized here as \textbf{RCWM-001: Rineya
Consequence World Model}. Earlier working names such as
reward-world-model and reward-token Transformer refer to this same
research line, not additional projects. {[}I03{]}

The research objective is not to demonstrate that numerical rewards can
be embedded. It is to demonstrate that \textbf{learning and reusing
consequence structure gives a better accuracy--cost--adaptation
trade-off than a comparably powerful scalar evaluator}. The most
persuasive result would be better action selection on new goals or
unseen compositions while holding environment experience, model
capacity, candidate plans and evaluation compute constant.

\subsection{The first decisive
experiment}\label{the-first-decisive-experiment}

Use a deterministic text/symbolic environment with exact state
snapshots, executable actions and an independent outcome checker.
Collect multiple action branches from the same starting state. Freeze
the Rineya language backbone. Compare a matched scalar Transformer, a
measured-feature predictor and a learned-consequence predictor. Hold out
complete problem families and goal combinations. Ask whether the learned
predictor chooses better actions using fewer new reward labels, not
whether its embedding plot looks attractive.

A successful first milestone is deliberately narrow: \textbf{on held-out
parents, the consequence representation supports better goal-conditioned
decisions at a measured cost and without changing the frozen language
model.} A null result is informative. It may show that the task lacks
reward diversity, the representation is insufficient, the proposed
actions are poor, or scalar prediction is already adequate.

\subsection{What this program is, and is
not}\label{what-this-program-is-and-is-not}

This is a research plan with an executable synthetic reference scaffold,
not a delivered trained Rineya world model. No production checkpoint is
included. No new Rineya capability gain, native GPU qualification,
remote branch, paid job or external account action was performed for
this bundle. The reference code checks selected mechanics in a small CPU
environment. It does not implement all later research arms.

World-model learning and reinforcement learning are overlapping
components, not competing taxonomies. A model of immediate rewards is a
reward model; a model of policy-conditioned cumulative rewards is a
value model. A predictor becomes a useful world model for this project
when its action-conditioned consequences correspond to independently
observable outcomes and improve control. Modeling Rineya's hidden-state
changes alone is an internal computation model, which must be evaluated
separately from an external environment model. Dreamer, successor
features and trajectory modeling provide precedents, but do not
establish that this Rineya adaptation will succeed. {[}S01, S02, S07{]}

\subsection{Five decisions to freeze before expensive
work}\label{five-decisions-to-freeze-before-expensive-work}

\begin{enumerate}
\def\labelenumi{\arabic{enumi}.}
\tightlist
\item
  \textbf{Subject:} the exact selected Rineya checkpoint, tokenizer,
  loader and executor versions; no substitution of a stock Qwen model or
  the older 344M student.
\item
  \textbf{Authority:} real environment feedback or an independent
  verifier defines correctness; predicted utility never grades its own
  success.
\item
  \textbf{Comparison:} the scalar evaluator receives the same context,
  action sequence and compute allowance as the consequence evaluator.
\item
  \textbf{Information:} predictions never receive future observations,
  actual future rewards or future actions chosen using hidden
  observations, unless those actions were explicitly fixed as an
  intervention.
\item
  \textbf{Promotion:} practical-effect, uncertainty, cost and retention
  gates are written before sealed evaluation.
\end{enumerate}

The first three reading priorities are successor features for the
factorization, DiSPO for distributions of achievable futures, and SF-GEN
for the direct language-generation precedent. Trajectory Transformer and
masked trajectory models establish that token-like trajectory
representations and flexible inference masks are already viable design
patterns. The Rineya contribution must therefore be a demonstrated
transfer, composition, efficiency or preservation improvement---not a
claim to have invented reward tokens. {[}S02, S04, S05, S07, S09{]}

\section{Scope, provenance and evidence
discipline}\label{scope-provenance-and-evidence-discipline}

\subsection{Rineya starting point}\label{rineya-starting-point}

The September 30 research handbook identifies a historical experimental
subject named \texttt{Rineya-Titan-Experimental-1.0.0-devfreeze}, with
596,049,920 base parameters. It also records that a laptop route passed
short native mechanics checks but encountered a prefix-versus-bulk
likelihood discrepancy during capture qualification. Those are
historical report statements, not a fresh audit of the current
repository. Stage G0 must re-establish the active checkpoint and current
qualification state before collecting hidden representations. {[}I01{]}

The target deployment envelope is the user's RTX 4070 Laptop GPU with 8
GB of VRAM and constrained host resources. This report uses that
envelope for design estimates; actual free memory, operating-system
overhead, thermal behavior and current RAM configuration must be
measured locally. Do not convert a short-sequence load success into a
claim that long-context training fits.

Preserve existing Rineya acquisition and retention work as independent
instruments. A base model can remain byte-identical while an attached
planner degrades deployed behavior. Consequently, report both
\textbf{weight preservation} and \textbf{system-level behavioral
preservation}. Existing composable-skill reports specifically motivate
paired evaluation, sealed splits and full resource accounting. {[}I02{]}

\subsection{Evidence labels}\label{evidence-labels}

\textbf{Literature finding} means a claim supported by a named primary
publication in its own setting. \textbf{Historical internal record}
means a statement from a previous Rineya report, not independently
reproduced here. \textbf{Derived example} means transparent mathematics
or a constructed example. \textbf{Proposal} means a design, threshold or
hypothesis requiring execution. \textbf{Reference mechanics} means code
actually run in the supplied synthetic CPU scaffold. \textbf{Native
result} is reserved for a measured experiment using the pinned Rineya
checkpoint; none is claimed in this delivery.

The annotated register lists 43 primary papers, author project pages and
official documentation references. Selected methods were examined for
Dreamer, successor features and DiSPO; many surrounding papers were
reviewed at the primary abstract or project-summary level. This is a
targeted synthesis, not an exhaustive systematic review or a full audit
of every cited codebase. Publication venue and version are separated
from access date. Recent preprints are explicitly not treated as
independent replication.

\subsection{Research authorization
boundaries}\label{research-authorization-boundaries}

This bundle authorizes no background job by itself. A later
implementation handoff should state the repository, parent revision,
permitted directories, model identity, data scope, compute budget and
stopping rule. The proposed branch is
\texttt{research/rineya-consequence-world-model-v1}; it has not been
created by this deliverable. Synthetic tests may run without a
checkpoint. Native extraction, training, model downloads, paid APIs,
cloud rental and remote pushes need the user's applicable authorization
and the repository's current policies.

The main scientific question should remain stable even if a particular
implementation fails. Preserve reports and failing tests instead of
deleting an experiment because it did not support the hypothesis.
Conversely, architectural elegance is not a reason to continue spending
compute after a well-controlled simpler baseline dominates.

\section{Formal model: actions, consequences and
goals}\label{formal-model-actions-consequences-and-goals}

\subsection{Define the state that the model actually
has}\label{define-the-state-that-the-model-actually-has}

Let \(o_t\) be an observation and \(H_t=(o_0,a_0,o_1,\ldots,o_t)\) the
available interaction history. The model receives a belief
representation \(b_t=E_\eta(H_t,e)\), where \(e\) contains environment
rules and available tools. This is not necessarily the true Markov
state. In partially observed tasks, two identical current observations
can imply different consequences because their histories differ.

An action is not merely the string \texttt{VERIFY}. Define

\[
a_t=(\text{operator},\text{arguments},\text{budget},\text{executor version}).
\]

The transition record includes the resulting observation, executed
status, duration, resource use and independent checks. Let
\(\phi_t=\Phi(\tau_t)\) encode the consequences of transition
\(\tau_t=(H_t,a_t,o_{t+1})\). A useful starting decomposition is

\[
\phi_t=[\phi_t^{\mathrm{measured}};z_t^{\mathrm{learned}}].
\]

Measured components anchor the model to checkable events. Learned
components preserve the user's proposed latent semantic space. Neither
is assumed to be intrinsically good or bad. A failed action, information
acquisition, elapsed computation and final task success are different
events and should remain distinguishable.

\subsection{Three related prediction
targets}\label{three-related-prediction-targets}

A one-step model predicts
\(p_\theta(\phi_t,b_{t+1},d_t\mid H_t,a_t,e)\), where \(d_t\) describes
termination. A plan-conditioned model predicts consequences for a
\textbf{specified} candidate action plan \(A\):

\[
p_\theta(\phi_{t:t+K-1},b_{t+1:t+K},d_{t:t+K-1}\mid H_t,A,e).
\]

A successor-feature model predicts a discounted accumulation under an
explicitly named continuation policy \(\pi\):

\[
\psi_K^\pi(b_t,a_t)=\mathbb{E}_{\pi}\left[\sum_{k=0}^{K-1}\gamma^k\phi_{t+k}\mid b_t,a_t\right].
\]

These targets are not interchangeable. The same first action can be
excellent under a competent continuation policy and poor under one that
repeatedly makes mistakes. Every successor target must therefore carry a
policy identity or the actual open-loop plan. A target mixed across
behavior policies without that information is a mixture, not an
optimal-action value.

For variable-duration actions, either use explicit per-decision
discounting or use elapsed-time discounting with cumulative duration.
State which convention is being estimated. A long tool call must not
quietly receive the same latency cost as a cheap token action. The
reference scaffold implements finite-horizon discounted features, not an
infinite-horizon policy-learning algorithm.

\subsection{Utility is a readout, not the world
itself}\label{utility-is-a-readout-not-the-world-itself}

A tractable first reward family is

\[
r_g(\tau_t)=w_g^\top\phi_t,\qquad Q_g^\pi(b,a)=w_g^\top\psi^\pi(b,a).
\]

This factorization is grounded in successor-feature research. Its
advantage is limited but useful: a new weighting of already represented
consequences can change values without relearning dynamics. It does not
mean that every possible new objective lies in the learned feature span,
or that a policy optimized for an old goal remains optimal for a new
one. Generalized policy improvement uses a set of supported policies;
planning can search other supported candidates. Neither guarantees
recovery of unobserved skills. {[}S02, S03{]}

A nonlinear utility \(u_g(\phi)\) is legitimate, but it changes the
requirements. In general,

\[
\mathbb{E}[u_g(\phi)]\ne u_g(\mathbb{E}[\phi]).
\]

A mean embedding that combines ``safe success'' and ``dangerous
failure'' is not an actual intermediate outcome. Risk constraints,
threshold rewards and conjunctions may require joint samples, richer
statistics or a distribution over successor features. DiSPO is a
particularly relevant precedent for distributions of achievable long-run
features; its dataset-supported reachability is an important limitation
to retain. {[}S04{]}

\subsection{Preference-only goals versus goals that alter
behavior}\label{preference-only-goals-versus-goals-that-alter-behavior}

When a goal simply says ``value speed more'' and the proposed plan is
fixed, the predicted physical consequences should not change. Only
utility should change. This is a valuable invariance test.

When the goal changes an action's actual execution---for example, a
\texttt{WRITE\_REPORT} operator whose prompt is goal-dependent---the
consequences can change. Supply the operative goal to the executor and
model. Do not enforce goal invariance where it contradicts the
environment. Distinguish \textbf{preference context}, \textbf{task
specification} and \textbf{control instruction} in the schema.

\subsection{What an embedding can and cannot
mean}\label{what-an-embedding-can-and-cannot-mean}

An embedding is a learned coordinate system. Its dimensions do not
automatically correspond to correctness, information or reasoning
phases. For an invertible matrix \(A\), setting \(\phi'=A\phi\) and
\(w'=A^{-\top}w\) preserves the scalar reward. Consequently, many
geometries can explain the same observations. Nearest-neighbor plots and
isolated cosine values are diagnostics, not proof of a unique semantic
decomposition.

Meaningful tests concern preserved information and behavior: can a small
readout recover held-out observable properties, can new goals be fitted
with fewer labels, do controlled action changes cause the predicted
consequence changes, and do those predictions improve action selection?
A cluster name is assigned after a blinded audit; it is not a training
truth because the cluster looks intuitively coherent.

There is also no preference-identification miracle in unsupervised data.
The same transition dataset is compatible with \(w\) and \(-w\).
Unsupervised learning can discover structure; something must ground
which outcomes are desirable. A verifier, task success flag, preference
comparison, explicit goal or chosen intrinsic objective supplies that
grounding. DIAYN and preference-learning work illuminate different sides
of this distinction. {[}S17, S20{]}

\section{Relationship to prior art and the defensible
contribution}\label{relationship-to-prior-art-and-the-defensible-contribution}

\subsection{Four families to
distinguish}\label{four-families-to-distinguish}

\textbf{Sequence-modeling policies.} Decision Transformer predicts
behavior from history and desired return. It is a useful baseline for
goal-conditioned action generation, but a desired return is not an
observed future outcome. Generalized Decision Transformer broadens the
conditioning statistics; it also highlights why hindsight targets and
deployment inputs must be separated. {[}S06, S18{]}

\textbf{Trajectory world models.} Trajectory Transformer gives the
closest early precedent for representing states, actions and rewards as
token sequences. UniMASK and Masked Trajectory Models broaden the tasks
expressible by masking. They motivate the interface, not a conclusion
that a causal language model should simply be switched to bidirectional
attention without requalification. {[}S07--S09{]}

\textbf{Task-relevant latent models.} Dreamer and MuZero show different
ways of making predictive models useful for decisions. DeepMDP and
bisimulation-oriented representation work motivate preserving
decision-relevant transition structure rather than reconstructing all
raw observations. These approaches also caution against an excessively
compressed representation that removes distinctions required by later
goals. {[}S01, S11--S13{]}

\textbf{Reusable future-consequence representations.} Successor features
separate consequences from a reward readout. DiSPO models richer
future-feature distributions. SF-GEN already brings successor features
to controlled language generation, so ``successor features for an LLM''
cannot be our novelty claim. {[}S02--S05{]}

\subsection{The research space that remains worth
testing}\label{the-research-space-that-remains-worth-testing}

The proposed contribution is a controlled combination: pretrained
compact-language representations; executable, parameterized reasoning or
tool actions; learned consequence geometry; goal revaluation;
long-distance and cross-subtask dependencies; and preservation-aware
integration. A combination is not scientifically important merely
because its parts can be connected. It becomes interesting when it
provides a reproducible benefit that a matched alternative cannot
explain.

Candidate publishable findings include a lower reward-label requirement
for new goals; improved choice among plans requiring held-out skill
compositions; improved accuracy at equal total reasoning cost; or a
learned representation that preserves useful consequence distinctions
across domains while scalar evaluators fail. An equally useful negative
finding would show that clustering adds interpretability but no control
value, or that long-context improvements arise entirely from additional
context rather than consequence factorization.

The comparison should not pit a full trajectory Transformer against a
scalar MLP with only one local state. A scalar model can use attention,
long context, rich hidden features and multiple outputs. Match those
capabilities before attributing a gain to the representation. This is an
experimental design requirement, not a theorem that the richer model
must win.

\subsection{Closest precedents: adoption
decisions}\label{closest-precedents-adoption-decisions}

\begin{longtable}[]{@{}
  >{\raggedright\arraybackslash}p{(\columnwidth - 4\tabcolsep) * \real{0.3333}}
  >{\raggedright\arraybackslash}p{(\columnwidth - 4\tabcolsep) * \real{0.3333}}
  >{\raggedright\arraybackslash}p{(\columnwidth - 4\tabcolsep) * \real{0.3333}}@{}}
\toprule\noalign{}
\begin{minipage}[b]{\linewidth}\raggedright
Reference
\end{minipage} & \begin{minipage}[b]{\linewidth}\raggedright
Adopt first
\end{minipage} & \begin{minipage}[b]{\linewidth}\raggedright
Defer or reject as an initial dependency
\end{minipage} \\
\midrule\noalign{}
\endhead
\bottomrule\noalign{}
\endlastfoot
Successor features & Explicit reward factorization and policy identity &
Claims of arbitrary zero-shot optimality \\
DiSPO & Direct future-feature prediction comparison & Full diffusion
outcome generator before a baseline works \\
SF-GEN & Frozen-LM feature/value comparator & Claims that this is the
first language application \\
Trajectory and masked models & Typed slots and audited masks &
Unrestricted future visibility during planning \\
Dreamer / MBPO & Re-anchoring and model-error gates & Long imagined-RL
rollouts at project start \\
CPC / VICReg / VQ & A staged representation ablation & All objectives
and codebooks enabled at once \\
\end{longtable}

The current literature register also includes text-world-model and
visual predictive-model work through 2026. Those sources refine the
roadmap but do not establish a Rineya result. Earlier conversation
percentages and venue attributions that were not reverified are not used
as evidence in this report. {[}S36--S39{]}

\section{Architecture and implementation
boundaries}\label{architecture-and-implementation-boundaries}

\subsection{One conceptual model, two implementation
routes}\label{one-conceptual-model-two-implementation-routes}

The intended architecture is a Transformer that relates history, actions
and future consequences in a shared computation. The safest first
experiment uses a \textbf{sidecar} so representation learning can be
debugged without changing Rineya. That is a staging choice, not
abandonment of the user's unified-model idea. A later \textbf{native
typed-slot adapter} implements the same contracts inside the actual
Rineya backbone. Compare both rather than assuming that sharing weights
is always more efficient.

For the sidecar, a frozen capture adapter extracts a bounded set of
Rineya context vectors. A small typed Transformer receives those vectors
and candidate actions. It outputs consequence features, next-state
summaries and optional uncertainty estimates. The goal readout scores
predicted consequences. A controller chooses a feasible action, executes
it in the real environment and re-encodes the actual new observation.

For the native route, introduce typed embeddings or projected continuous
slots at Rineya's verified hidden width. Train limited adapters and
small output heads while retaining the ordinary language route. Use the
supported embedding-input interface or a tested wrapper; do not silently
resize the language vocabulary and retie its output weights without
checking exact language-logit preservation. Stock Qwen3 documentation is
an interface reference, not proof that a custom Rineya loader implements
it identically. {[}S42{]}

\subsection{Proposed component
contracts}\label{proposed-component-contracts}

\begin{longtable}[]{@{}
  >{\raggedright\arraybackslash}p{(\columnwidth - 4\tabcolsep) * \real{0.3333}}
  >{\raggedright\arraybackslash}p{(\columnwidth - 4\tabcolsep) * \real{0.3333}}
  >{\raggedright\arraybackslash}p{(\columnwidth - 4\tabcolsep) * \real{0.3333}}@{}}
\toprule\noalign{}
\begin{minipage}[b]{\linewidth}\raggedright
Component
\end{minipage} & \begin{minipage}[b]{\linewidth}\raggedright
Input
\end{minipage} & \begin{minipage}[b]{\linewidth}\raggedright
Output / invariant
\end{minipage} \\
\midrule\noalign{}
\endhead
\bottomrule\noalign{}
\endlastfoot
History encoder & Available text, evidence and events & Time-stamped
context slots; no future facts \\
Action encoder & Operator, arguments, budgets, version & Typed action
vectors; executable semantics \\
Consequence predictor & History and permitted action prefix/plan &
Consequence distributions or embeddings \\
Outcome decoder & Consequence representation & Independent observable
targets for grounding \\
Goal readout & Goal weights or learned task embedding & Utility; no
alteration of physical facts \\
Candidate proposer & Available state and legal operator set & Bounded
candidate plans with validity checks \\
Controller & Predictions, support, constraints & One admitted action or
a safe fallback \\
Recorder & Actual action and actual result & Immutable
training/evaluation evidence \\
\end{longtable}

The predictor must not receive its targets in its forward signature. In
particular, a tensor named \texttt{future\_states} should not be
accepted by the prospective evaluator. Retrospective representation
training may encode outcomes to create targets, but target construction
lives in a separate module with a separately audited visibility
contract.

\subsection{Typed token construction}\label{typed-token-construction}

For event \(i\), form a projected vector

\[
x_i=P_{\mathrm{type}(i)}v_i+e_{\mathrm{type}(i)}+e_{\mathrm{time}(i)}+e_{\mathrm{subtask}(i)}.
\]

Text uses the ordinary tokenizer and embedding path. Continuous state
summaries and consequence slots use learned projections. Operator
embeddings distinguish operations; argument encodings distinguish
\texttt{CALCULATE(2+2)} from \texttt{CALCULATE(large\_expression)}.
Duration and resource encodings prevent the model from treating an
expensive verification process as a free semantic label. Type IDs
distinguish a proposed outcome from an observed outcome even when their
dimensions match.

A candidate serial sequence can be serialized as history slots, an
action slot, an outcome-query slot, another action slot and another
query. Query slots contain learned placeholders, not ground-truth
outcomes. In the first causal version, the query for step \(t\) can
access only the initial observed history and candidate actions through
\(t\). The model learns to integrate the consequences of earlier planned
actions implicitly; it does not receive teacher-forced future
observations at evaluation.

Outcome vectors do not have to be literal vocabulary tokens. Separate
projections and heads can preserve the useful Transformer interface
without allocating thousands of arbitrary reward strings. Learned
discrete codes are a later representation choice, not a prerequisite for
the architecture.

\subsection{Pilot tensor
specification}\label{pilot-tensor-specification}

The following are proposed screen values, not measured optima. Start
with a 128--256-dimensional sidecar, 2--4 layers, 4 attention heads,
32--64 learned consequence dimensions and a small set of grounded
features. For native integration, infer width and layer layout from the
checkpoint. Do not hard-code stock-model dimensions into persisted
captures.

A training batch contains context \texttt{{[}B,Lc,dc{]}}, proposed
actions \texttt{{[}B,H{]}} plus argument features, target consequences
\texttt{{[}B,H,F{]}}, target-validity masks \texttt{{[}B,H,F{]}},
step-validity masks \texttt{{[}B,H{]}}, optional next-observation
targets and a split/parent index. Goal weights may have shape
\texttt{{[}B,F{]}} for a measured-feature baseline; learned-latent
readouts must fit weights in that latent space using allowed calibration
data. A hand-written six-dimensional goal vector cannot simply be dotted
into an arbitrary learned 64-dimensional latent.

The predictor should expose both \texttt{predict\_step} and
\texttt{predict\_plan}. If these share weights, specify their masks.
Output a mean and per-target uncertainty only where the training
likelihood makes this meaningful. A vector of independent standard
deviations is not a full covariance model. Store units and
feature-normalization metadata with every model artifact.

\subsection{Backbone reuse and target
stability}\label{backbone-reuse-and-target-stability}

Extract representations at explicit action boundaries. Compare
last-token, pooled evidence and multiple selected-layer summaries on
development data. Normalization and pooling must be consistent across
capture modes. Do not assume a raw subtraction of hidden states from
different prompts measures a meaningful semantic transition.

A frozen encoder gives stable targets and cheap caching. A learned
encoder can improve task relevance, but then the predictor chases moving
targets. Use a separately versioned target encoder, stop-gradient
targets and a controlled exponential moving average only after the
frozen-target baseline works. Recompute or invalidate caches whenever
encoder weights, tokenizer, prompt template, mask, precision or capture
location changes. A stale representation cache is a scientific data
error, not merely an engineering inconvenience.

The CPU scaffold intentionally uses a tiny synthetic state input rather
than pretending to load Rineya. Its learned latent bottleneck is
grounded through feature prediction. It demonstrates the interface and
tests selected invariants; the fully self-supervised latent-consequence
arm and native checkpoint integration remain planned work.

\section{Executable actions and consequence
semantics}\label{executable-actions-and-consequence-semantics}

\subsection{Give each token an operational
definition}\label{give-each-token-an-operational-definition}

\texttt{RETRIEVE} must name a retrieval corpus, query generator,
top-\(k\), evidence budget and tool version. \texttt{VERIFY} must
identify what is checked and the verifier's authority. \texttt{RECUR}
must specify a qualified recurrent module, depth, gating value and
termination rule. \texttt{ANSWER} must identify the format and final
evaluator. Two runs with the same operator name but different prompts or
tool versions are different action mechanisms.

An action record contains its ID, parent events, arguments,
preconditions, anticipated resource needs and permitted side effects.
The executor returns a status from a fixed set such as success, invalid
precondition, timeout, tool error or explicit refusal. A failure is not
silently dropped from training. Dropping unsuccessful calls biases the
predictor toward imaginary reliability.

For the first real Rineya pilot, use only a few controllable operations:
retrieve a document span from a fixed corpus, execute a bounded
arithmetic expression, check an answer against a symbolic constraint,
and submit a final answer. Broader cognitive labels such as
``decompose'' or ``think harder'' can follow once their execution
contracts are stable. Learned action representations are useful
precisely when there is actual structure in the available operations.
{[}S10{]}

\subsection{Grounded feature
dictionary}\label{grounded-feature-dictionary}

Use a versioned dictionary with units, observability and label source. A
initial example includes terminal task success, a newly verified
constraint, contradiction detected, useful evidence obtained, elapsed
milliseconds, generated tokens, tool failure and irreversible constraint
violation. Some are binary, some counts, some nonnegative costs. Predict
their distributions or normalized continuous values as appropriate.

Do not equate an observed reduction in the model's entropy with a
correct conclusion. Call it ``entropy change,'' not ``information
gain,'' unless the experiment provides a well-defined information
target. Similarly, ``progress'' requires an external measure such as
satisfied prerequisites or a reduction in known distance to a simulator
goal. A fluent narrative that says progress occurred is not a
measurement.

Earlier brainstorming placed retention beside compute and correctness in
a reward vector. For ordinary inference actions, that is misleading: the
frozen model does not change its weights when it retrieves a document.
Retention belongs in \textbf{cross-update evaluation}, unless the
selected action is itself a training update with measured downstream
effects. Keep the feature inventory honest about the timescale of each
quantity.

\subsection{Avoid double-counted
reward}\label{avoid-double-counted-reward}

A step that verifies a prerequisite can deserve a shaping signal, but
the final success reward should not be paid repeatedly after success.
The plan objective must specify whether intermediate signals are
explanatory features, auxiliary supervision or actual additive rewards.
Those roles differ.

For example, store elapsed time as a positive measured quantity and
place a negative coefficient in the utility readout. Do not encode
negative time and also multiply by a negative penalty. Count
candidate-generation and scoring costs in system reporting even though
they do not appear in the environment transition. An improved controller
that spends ten times as much on evaluation has not demonstrated compute
efficiency.

\subsection{A constructed example}\label{a-constructed-example}

A task requires a key, an opened door and a collected object. A shortcut
collects the object immediately but causes damage. Consequence features
are \texttt{{[}object\_gained,\ damage,\ action\_cost{]}}. A safe
three-step route yields a cumulative vector \texttt{{[}1,0,3{]}}; a
one-step shortcut yields \texttt{{[}1,1,1{]}}.

With weights \(w_{\mathrm{safe}}=[10,-20,-1]\), the route scores \(7\)
and the shortcut scores \(-11\). With \(w_{\mathrm{fast}}=[10,-1,-2]\),
the route scores \(4\) and the shortcut scores \(7\). These are
constructed arithmetic values, not model predictions or empirical
results. The physical futures did not change; the preference did.

This example is intentionally simple. If success also requires no
damage, then utility is a conjunction rather than this linear readout.
Add the relevant joint feature, predict a joint distribution, or use a
nonlinear scorer with an appropriate outcome distribution. Do not use
linear scoring outside the feature family and blame the world model for
an incorrectly specified objective.

\section{Learning the consequence space and testing
clustering}\label{learning-the-consequence-space-and-testing-clustering}

\subsection{Preserve the discovery
hypothesis}\label{preserve-the-discovery-hypothesis}

The first measured-feature model is a debugging instrument. The actual
research hypothesis is that a representation learned from transitions
contains useful distinctions we did not fully prescribe. A credible
program therefore needs three parallel representation tracks: measured
features, learned continuous features, and learned discrete or clustered
features. Each must be compared using the same action data and the same
downstream reward-label budget.

For learned continuous features, encode observed transition or
trajectory outcomes with a target encoder and predict that
representation from the preceding history plus action. Add a small
decoder for independently checkable properties. This makes it possible
to detect a representation that predicts itself perfectly while carrying
no useful information. Use a held-out linear probe and a nonlinear probe
separately: a high-capacity probe can hide a poorly organized
representation.

Predictive contrastive learning, variance/covariance regularization and
vector quantization supply alternative mechanisms. They should be
introduced one at a time. None guarantees that the resulting latent
dimensions correspond to useful rewards. {[}S14--S16{]}

\subsection{Clustering protocol}\label{clustering-protocol}

Begin with train-only standardization and a continuous baseline. Fit
\(k\)-means to TRAIN embeddings for a small predeclared set such as
\(K\in\{16,64,128\}\). Freeze centers before development or test
assignment. Compare hard assignments, soft distances and the unclustered
vector under equal downstream model capacity. A random projection and
random-label cluster control help identify whether a gain comes from
regularization or parameter count rather than semantic grouping.

For a vector-quantized arm, log code usage, entropy, empty/dead codes,
assignment turnover, quantization distortion and downstream regret.
Resetting dead codes is a training intervention that must be logged and
must not use evaluation examples. A codebook can stabilize a vocabulary
of outcomes, but can also merge rare catastrophic failures with frequent
harmless events. Evaluate failure distinctions explicitly rather than
relying on average reconstruction error.

Inspect clusters only after the quantitative screen. Draw a blinded
sample per cluster, have an evaluator describe observable commonalities,
then test those descriptions on fresh examples. Do not let a compelling
hand-assigned name turn into an unsupported assertion that the model has
discovered a cognitive primitive.

\subsection{Sparse grounding
experiment}\label{sparse-grounding-experiment}

Use nested label budgets, for example 1\%, 5\%, 10\%, 25\% and 100\% of
the same TRAIN parent set. All representation arms receive the same
unlabeled transitions; all receive identical labeled parents at each
budget. Count verifier calls, human comparisons, teacher generation,
rejected traces and environment execution separately. ``Unlabeled''
transitions may still be expensive to collect.

Fit the utility readout using only the permitted calibration labels.
Separate two transfer tests. \textbf{In-span revaluation} changes
weights on already represented features. \textbf{New-outcome transfer}
introduces a property not explicitly included in prior grounding. The
latter can benefit from learned latent structure, but requires evidence
that the representation retained that property. Neither test may reuse
variants of held-out parents in pretraining or clustering.

The principal label-efficiency endpoint is the number of grounding
labels required to reach a preregistered action-regret or success
target. Report uncertainty on that curve and the total data-production
cost. A claim that clustering uses ``10\% of labels'' is not justified
by selecting one favorable random label subset.

\subsection{Representation failure
diagnostics}\label{representation-failure-diagnostics}

Track feature variance, covariance spectrum, nearest-neighbor stability,
decoder quality, cross-goal sufficiency, action sensitivity and nuisance
sensitivity. Compare prediction under real actions against shuffled
action labels. If shuffling barely matters, the model may predict
outcome frequencies from the state while ignoring the proposed
intervention.

Do not force all successful outcomes to be neighbors. Different
successful routes can have different costs, risks and prerequisite
effects. A contrastive objective that collapses them may improve a
binary success probe while harming goal transfer. Likewise, ``different
action'' is not automatically a negative pair: two tools may cause the
same observable effect.

\section{Information flow, masking and counterfactual
validity}\label{information-flow-masking-and-counterfactual-validity}

\subsection{Three masks, three
meanings}\label{three-masks-three-meanings}

\textbf{Prospective causal transition mask.} A query for the consequence
of action \(a_t\) may see the observed prefix and candidate actions
through \(t\). It may not see later actions, outcomes or rewards.
Changing \(a_{t+1}\) must not change the predicted consequence at \(t\).
This is the default dynamics interface.

\textbf{Prospective whole-plan mask.} A plan-value or plan-outcome query
may see a fully specified open-loop candidate plan, including later
proposed actions. All unobserved future outcomes remain absent or
masked. This can score the total plan, but it should not claim that
later choices physically cause earlier consequences. A suffix-aware
auxiliary scorer and a causal dynamics head may share some parameters
while maintaining separate visibility tests.

\textbf{Retrospective mask.} A representation learner may encode a
completed trajectory to learn outcome structure or redistribute credit.
Its output is a training target or retrospective diagnostic, not a
prospectively available state. Confusing this mode with the deployed
evaluator is a particularly easy way to manufacture excellent validation
scores. Masked trajectory research motivates these interfaces;
deployment validity is our responsibility. {[}S08, S09, S18{]}

\subsection{The subtle future-action
leak}\label{the-subtle-future-action-leak}

Even when rewards and observations are masked, a recorded future action
can reveal a hidden observation. A logged policy may choose
\texttt{GO\_LEFT} only after seeing a secret clue. Showing that future
action to an earlier predictor leaks the clue. An action column is not
safe merely because it is called an action.

For the first plan-conditioned dataset, choose candidate plans
\textbf{before} executing their branches, using only the shared starting
observation and allowed information. Record
\texttt{plan\_origin=open\_loop\_precommitted} or a policy-tree contract
with explicit observable branch conditions. For retrospective
trajectories selected by an adaptive policy, use causal prefixes or
include an honest policy-conditioned target; do not treat the realized
suffix as a freely chosen plan.

\subsection{Predictive association versus causal
effect}\label{predictive-association-versus-causal-effect}

Logged trajectories can be confounded by environment difficulty, hidden
state and behavior-policy competence. To learn an action effect, prefer
exact simulator snapshots and branching interventions: restore the same
starting state, execute different permitted actions, and compare actual
consequences. Keep all branches in the same data split. In stochastic
tasks, log random seeds and use repeated draws; shared random numbers
can reduce variance but do not turn stochastic outcomes into
deterministic facts.

A language-model hidden-state intervention is not automatically an
intervention on an external task variable. First test whether its
predicted changes correspond to observable environment changes. Only
then study stronger internal causal hypotheses. Attention maps, cluster
assignments and decoded labels provide hypotheses for interventions, not
proof of causal structure.

\subsection{Tests that must fail before the fix and pass after
it}\label{tests-that-must-fail-before-the-fix-and-pass-after-it}

Perturb a hidden future reward and confirm the prospective input bytes
and outputs are unchanged. Perturb future action suffixes and confirm
causal earlier outputs are unchanged. Introduce a future timestamp in a
supposedly known observation and reject the record. Replace a legal
precommitted plan with a behavior-selected suffix and require a
different mode or an error. Change padding, batch order and unrelated
examples and verify the expected numerical tolerance.

Mask Boolean conventions deserve explicit tests: the
\texttt{torch.nn.Transformer} API marks blocked positions with
\texttt{True}; other attention APIs can use the opposite convention. A
shape-correct mask can therefore be completely wrong. Pin the installed
API and test information flow behavior, not just matrix dimensions.
{[}S41{]}

\subsection{Parallel output is not joint
uncertainty}\label{parallel-output-is-not-joint-uncertainty}

A Transformer can produce reward or feature predictions at every query
position in one pass. Independent heads at those positions generally
estimate marginals. They do not guarantee coherent joint samples across
steps, tasks or reward dimensions. If a task requires ``both subtasks
succeed'' or a correlated failure mode, test the joint event directly.
Later options include a shared latent scenario variable, an
autoregressive consequence decoder or a distributional trajectory model.
Count their sampling cost in comparisons.

\section{Data program and environment
ladder}\label{data-program-and-environment-ladder}

\subsection{Start where errors can be
localized}\label{start-where-errors-can-be-localized}

The first environment should be simple enough that an oracle can
exhaustively evaluate candidate plans. The supplied KeyDoor simulator is
a mechanics example, not the final benchmark. Build a broader symbolic
suite with prerequisite chains, bounded program execution,
document-evidence retrieval and resource-constrained parallel subtasks.
Each environment should expose immutable snapshots, legal-action
validation, independent outcome checking and a controllable observation
function.

Use separate axes for context length, decision horizon, branching factor
and partial observability. A 16,000-token prompt with one action is not
a long-horizon task. A 20-action sequence with a tiny state is not a
long-context reading task. A fair evaluation varies one axis at a time,
then tests their combinations.

Only after the exact simulator gates should the program move to
TextWorld or ALFWorld-style interactions, and later to a sandboxed web
benchmark. These environments offer richer interactive structure but
introduce additional language and state-aliasing failure modes. External
benchmark acquisition, licenses and versions remain a preflight task. No
real purchases or account changes belong in the training environment.
{[}S31--S33{]}

\subsection{Collection policy and
coverage}\label{collection-policy-and-coverage}

Collect a mixture of legal random exploration, heuristic actions, the
frozen Rineya proposal policy and deliberately selected alternatives.
Log which policy selected each action and its available information.
Retain failures, invalid attempts, no-ops and timeouts with explicit
status labels. An all-successful dataset teaches the model little about
the actions it most needs to reject.

For a practical pilot, propose 2,000 independent training parents, 300
development parents and 500 sealed evaluation parents, with several
matched branches per parent. These are \textbf{starting budget
proposals}, not a power calculation or a guarantee of sufficient
coverage. Use the development pilot to estimate variance and effect
sizes; freeze the final sample size before opening the sealed set. If
the required sample size exceeds the compute budget, narrow the claim or
improve measurement rather than reporting an underpowered absence of
harm.

Balance training by informative transition types, not only by trajectory
count. A million repeated no-ops can dominate reconstruction loss.
Preserve natural prevalence for a separate calibration set so a balanced
training sampler does not produce misleading real-use probabilities.
Report support for each operator, argument type, context band, goal
family and failure mode.

\subsection{Data contract}\label{data-contract}

Every trajectory record needs a schema version, parent ID,
counterfactual family ID, episode ID, split, environment/version hash,
initial observation, source history and observation timestamps. Events
carry action/operator IDs, arguments, executor hash, preconditions,
start and end times, dependencies, actual outputs, statuses, feature
values, feature-validity masks, reward-source IDs, termination and
truncation. Optional hidden-state references carry
model/tokenizer/capture hashes and byte offsets.

Keep \textbf{observed facts}, \textbf{proposed actions},
\textbf{predicted futures} and \textbf{retrospective annotations} in
distinct fields. Do not overload an \texttt{outcome} field with
whichever version is currently convenient. A typed schema reduces the
chance that a training target later becomes an inference feature.

Prediction logs need candidate-set identity, all candidate scores,
uncertainty/support estimates, chosen action, fallback reason, actual
outcome and resource totals. Storing only the selected action hides
whether gains come from the proposer or scorer. Store all candidates
before execution so later analysis cannot silently change the comparison
set.

\subsection{Split integrity}\label{split-integrity}

Partition by a stable problem-family key before generating branches,
paraphrases, goal variants or labels. All descendants remain together.
When tasks share a generator template or source document, use a
sufficiently broad grouping key to prevent template leakage. A random
row split of trajectories from the same parent is not a generalization
test.

Normalizers, codebooks, representation pretraining, reward readouts and
hyperparameter selection must respect the split boundary. Separate a
\textbf{new-goal calibration set} from the held-out evaluation parents.
New goal labels may be permitted during adaptation, but labels from
evaluation scenarios are not. Record this distinction in the experiment
contract rather than calling every goal adaptation ``zero-shot.''

Preserve a frozen ledger of hashes for raw data, parent partitions and
generated variants. Deduplicate before split finalization and recheck
after generation. Where a benchmark has its own train/test partition,
retain it and add internal development splits only within training data.

\subsection{Label provenance and
privacy}\label{label-provenance-and-privacy}

Each feature label records how it was obtained: simulator state,
executable checker, deterministic tool, human comparison or learned
proxy. Report accuracy against independent checks for any learned
labeler. Use separate model identities for generator and evaluator when
appropriate, and disclose shared training or correlated errors. Never
label the reward model's own prediction as ground truth.

Do not include private email, real credentials, personal identifiers or
proprietary documents in synthetic training data by default. Tool
outputs are untrusted data, not instructions to alter the experiment.
Redact secrets before logging; retain enough provenance to reproduce a
sanitized example. Store only the hidden representations needed for the
experiment, with explicit access and retention policies.

\section{Training recipe and staged
objectives}\label{training-recipe-and-staged-objectives}

\subsection{Stage A: measured consequences as an
instrument}\label{stage-a-measured-consequences-as-an-instrument}

Train a small model to predict known consequence features from initial
state/history plus a precommitted action prefix. Use binary
cross-entropy for Bernoulli events, cross-entropy for categorical state
variables and a suitable regression likelihood for measured continuous
quantities. The CPU scaffold uses mean-squared error for simple
synthetic numerical features; production binary probabilities require
the corresponding calibrated heads.

Normalize continuous targets using TRAIN statistics. Preserve physical
units in the artifact. Apply feature-validity masks so an unavailable
label is not treated as zero. Start with a fixed-length short horizon
and a single frozen encoder. Require the model to overfit a tiny
diagnostic set before larger experiments; failure there usually
indicates a contract, optimization or observability problem.

\subsection{Stage B: learned consequence
representations}\label{stage-b-learned-consequence-representations}

Let \(z_t=E_{\bar\eta}(\tau_t)\) be a target outcome representation and
\(\hat z_t=F_\theta(H_t,a_t)\) its prediction. A basic predictive loss
is

\[
\mathcal L_{\mathrm{pred}}=\frac{1}{|M|}\sum_{t\in M}\|\hat z_t-\operatorname{sg}(z_t)\|_2^2.
\]

The stop-gradient operator prevents the predictor from trivially moving
both sides through that term, but does not alone prevent the whole
representation from collapsing. Combine it with a grounded decoder,
contrastive prediction or an explicitly tested variance/covariance
objective. Start with one anti-collapse mechanism, not all of them. CPC
and VICReg provide relevant objective families; our exact transition
augmentations and feature targets remain hypotheses. {[}S14, S16{]}

An auxiliary observable loss can decode state predicates, action status
and elapsed resources. Evaluate both the auxiliary prediction and
downstream utility. A representation might predict all easy surface
properties while dropping a rare critical prerequisite. Stratified
diagnostics and held-out conjunctions are needed to expose that failure.

\subsection{Stage C: sparse utility
grounding}\label{stage-c-sparse-utility-grounding}

Freeze or slowly adapt the consequence encoder and fit \(u_\omega(z,g)\)
using verified returns or preference pairs. Use a linear readout first
because it exposes what is actually reusable. Compare a small nonlinear
readout only after the linear result is known. For pairwise labels,
train a calibrated preference model rather than pretending that every
comparison reveals an absolute numeric reward.

When fitting goal weights, constrain them to the declared feature family
or record the extrapolation. New-language goal prompts must be
translated to a versioned task representation and tested against
explicit objective definitions. Do not claim immediate generalization
from a natural-language goal embedding without evaluating that mapping.

\subsection{Stage D: future features and
consistency}\label{stage-d-future-features-and-consistency}

Add direct predictions of discounted future consequences at horizons 1,
2, 4 and 8. For a policy-conditioned successor head, enforce the correct
continuation-policy identity. A finite-horizon consistency target has
the form

\[
\psi_K^\pi(b_t,a_t)\approx\phi_t+\gamma\,\mathbb E_{a'\sim\pi}[\psi_{K-1}^\pi(b_{t+1},a')].
\]

Bootstrap only where continuation semantics are valid. True terminal
states have no future contribution. A time-limit truncation is not
automatically a terminal state; use an explicit modeling convention. The
environment API distinguishes these cases, but the learner must still
handle them correctly. {[}S43{]}

Compare direct successor prediction with recursive one-step rollouts.
The former can reduce accumulated prediction steps but hides
intermediate structure; the latter supports detailed replanning but can
accumulate error. A hybrid is a research option, not the default until
these failure modes are measured. {[}S04, S22{]}

\subsection{Objective budget and
optimization}\label{objective-budget-and-optimization}

A later combined objective may contain predictive, observable, reward,
variance, contrastive, codebook and retention terms. Record each
coefficient, its units, its gradient contribution and its activation
stage. Do not present a single long loss equation as an implementation
recipe. Begin with observable prediction, then add one mechanism at a
time.

For the sidecar, propose AdamW with a small development-only sweep over
learning rates such as \(10^{-4}\), \(3\times10^{-4}\) and \(10^{-3}\),
fixed token/transition exposure, early stopping on an independent
development criterion, gradient clipping and finite-loss checks. These
are engineering starting points, not paper-derived optimal settings. Log
samples seen, optimizer steps and total modeled tokens. Match baselines
on training examples and total tuning budget, not only epochs.

For native adapter training, verify forward/gradient parity on a small
case, then train only registered parameter sets. LoRA and quantized-base
adaptation are reasonable mechanisms, but their efficiency does not
certify retention or numerical fidelity. Quantization must be evaluated
as a separate variable. {[}S27, S28{]}

\subsection{Representation drift and
replay}\label{representation-drift-and-replay}

When the representation encoder changes, old goal weights and codebook
centers can become invalid. Either keep the encoder frozen, version and
refit dependent readouts using allowed data, or train an explicit
alignment map and evaluate its error. Store representation version in
every prediction record. A resumed training run that silently combines
old caches with new encoder weights is invalid.

Retain a small, fixed rehearsal set for language behavior and previous
task capabilities only when native weights are adapted. Do not use
sealed retention tests as rehearsal. Record update order and test
acquisition/retention separately so a model cannot pass preservation
merely by learning nothing new.

\section{Planning, uncertainty and real-environment
feedback}\label{planning-uncertainty-and-real-environment-feedback}

\subsection{Separate proposing from
evaluating}\label{separate-proposing-from-evaluating}

Generate a bounded set of legal action plans using fixed heuristics or
the frozen Rineya proposer. Validate syntax, preconditions, dependency
constraints and permissions before scoring. Cache the candidate set and
its hash. Every evaluator in a comparison receives exactly that set.

For an admitted candidate \(A\), score an expected task utility and a
development-calibrated model-risk penalty:

\[
J(A,g)=\mathbb E[U_g(\widehat\Phi(A))]-\beta\,\widehat{\mathcal U}(A).
\]

This is a heuristic conservative score unless a formal bound is
separately established. Enforce hard constraints through a feasible-set
filter, not through arbitrarily large learned penalties. A permission
violation must not become acceptable because the model predicts enough
reward elsewhere.

Execute only the first admitted action, observe the actual result,
update the history and replan. An open-loop plan is a proposal, not a
commitment to ignore new evidence. A safe fallback can be a verified
tool, an inexpensive direct answer on a qualified class, a request for
missing information, or abstention under the task's rules.

\subsection{What uncertainty means}\label{what-uncertainty-means}

Separate stochastic outcome variability from lack of model knowledge. In
a deterministic toy environment, residual uncertainty primarily reflects
model error or missing state; in a stochastic environment, repeated
execution can vary even under a perfect model. Ensembles are a practical
baseline, not an oracle for ignorance. Test whether uncertainty predicts
error under held-out actions, longer horizons and shifted contexts.
{[}S24, S25{]}

Report calibration by operator and difficulty as well as overall. Use
Brier score or log loss for probabilities; reliability plots, empirical
coverage and interval sharpness for continuous outcomes. Define the
abstention threshold on development data and freeze it. Measure coverage
and real task quality together: rejecting everything is not a useful
uncertainty system.

Candidate search itself can exploit the evaluator. Sweep plan count and
horizon, and compare predicted gains with actual utility. If predicted
reward increases while real utility falls, the controller is exploiting
model error. Keep failed cases as a dedicated stress set rather than
retraining on sealed evaluation failures. Proxy-reward overoptimization
and offline distribution shift make this an essential test. {[}S23,
S26{]}

\subsection{Direct selection, distillation and
RL}\label{direct-selection-distillation-and-rl}

The first controller can choose the highest-scoring permitted action
without online policy-gradient training. That does not make it
fundamentally unrelated to RL: it is a learned decision system trained
from interaction data. The question is which learning components are
needed, not which label is more impressive.

A later distilled policy can imitate the choices of a qualified, more
expensive planner, reducing inference cost. Verify that the teacher's
preference remains valid under real outcomes. Imagined RL is later
still: mix short simulated trajectories with real anchors, use
conservative support checks and validate every promoted policy in the
actual environment. The Dreamer and MBPO inspirations support this
staged idea, not a reason to start with unrestricted self-generated
reward loops. {[}S01, S22{]}

\subsection{Evaluation of control versus
representation}\label{evaluation-of-control-versus-representation}

Keep three levels of results separate: prediction fidelity, decision
quality within a fixed candidate set, and end-to-end agent performance
including proposal failures. A strong scorer cannot select an action
that was never proposed. Report the candidate oracle gap to show whether
further improvement should target the proposer or evaluator.

A model can predict many irrelevant details accurately yet rank actions
poorly. Conversely, it can rank familiar candidates well without
generalizing dynamics. Both are legitimate limited results. The ``world
model'' label should be scoped to the evaluated environment, observable
variables and intervention coverage. Recent language-world-model work
similarly motivates evaluating fidelity and actual agent utility
separately. {[}S38{]}

\section{Long context, delayed credit and parallel task
structure}\label{long-context-delayed-credit-and-parallel-task-structure}

\subsection{Long-context evaluation
grid}\label{long-context-evaluation-grid}

Propose context bands of 512, 2,048, 4,096 and 8,192 input tokens,
subject to measured hardware limits; action horizons of 1, 2, 4 and 8;
and dependency distances varied independently. Hold total relevant
information fixed while adding distractors in one test. In another,
increase the number of genuinely relevant prerequisites. A context-size
configuration is not a claim that the model uses that context
effectively.

Use exact synthetic references for early tests: a constraint introduced
near the beginning changes which final action is legal, two distant
facts must be combined, or an intermediate tool response invalidates an
earlier assumption. Report position sensitivity and truncation behavior.
LongBench and RULER offer later external context-evaluation references,
but do not replace interactive tests. {[}S29, S30{]}

For long raw prompts, compare full-history capture, bounded evidence
retrieval and hierarchical event summaries. Count encoding and retrieval
costs. A compressed history can make the sidecar cheap while discarding
the only fact that distinguishes two actions. Test that possibility
using adversarially placed prerequisites rather than trusting a
summarizer's own explanation.

\subsection{Delayed credit without hindsight
leakage}\label{delayed-credit-without-hindsight-leakage}

A useful early action may have no immediate success reward. Supervise
the eventual consequences while preserving when evidence becomes
observable. A retrospective return-decomposition model can propose
credit assignments; a prospective decision model must predict those
outcomes from information available before acting. RUDDER is relevant to
delayed credit, but an attention link from an early action to a final
reward is not automatically causal attribution. {[}S19{]}

A useful ablation compares terminal-only grounding, measured
prerequisite features and learned future-consequence embeddings. Match
the final objective and label costs. Intermediate supervision may
improve training, but then the gain is not purely architectural; the
study should identify the contribution of extra feedback.

\subsection{Represent parallel tasks as events and
dependencies}\label{represent-parallel-tasks-as-events-and-dependencies}

A parallel task is not simply a serial action list with subtask tags.
Define a directed acyclic graph of prerequisite relationships. Each
event has a start time, completion time, dependencies, resource claims
and output-availability time. Two actions may be independent in logic
but compete for the same GPU, file lock or API quota. A plan is feasible
only when both dependency and resource constraints hold.

Begin with synchronous rounds: independent actions launch together, then
their outcomes become available together. Later add asynchronous
completion, failures and cancellation. A query at time \(t\) may attend
only to outcomes available by \(t\), not to a sibling task that finishes
later. Proposed action metadata can be visible according to the selected
planning mode; actual results cannot arrive early because their token
was serialized earlier.

Test concurrency benefits against a serial controller at equal resources
and against an oracle scheduler. Metrics include task success, makespan,
resource cost, constraint violations, wasted work and cross-subtask
error propagation. A shorter action sequence is not necessarily a
shorter wall-clock run.

\subsection{Cross-task consequence
interactions}\label{cross-task-consequence-interactions}

When two tasks must both succeed, marginal success probabilities are
insufficient unless independence is justified. A shared failure source
can make the joint success probability very different from the product
of marginals. Include joint outcome labels, shared scenario variables or
sampled joint trajectories for this stage.

Test whether the model learns reusable prerequisites rather than task
IDs. Hold out new compositions of familiar operations, permutations of
independent subtasks and changes in resource contention. Preserve
counterfactual families in a single split. A system that succeeds only
on the serialized order seen in training has not learned the claimed
parallel structure.

\subsection{What attention does not buy for
free}\label{what-attention-does-not-buy-for-free}

Attention provides a mechanism to connect positions and train on many
positions in parallel. It does not enumerate an exponential plan space,
make all future observations available, infer causal structure from any
dataset, or execute tools simultaneously. The plan must still have a
proposer, a feasibility check and a measured execution budget. The
project's hypothesis is that shared contextual computation improves
prediction and ranking---not that attention abolishes planning.
{[}S40{]}

\section{Experimental program and promotion
gates}\label{experimental-program-and-promotion-gates}

\subsection{Hypotheses that can be
falsified}\label{hypotheses-that-can-be-falsified}

\textbf{H1, useful representation:} learned consequence features improve
independently measured candidate selection over a matched scalar
Transformer. Falsify the claimed representation advantage when the
effect disappears after matching context, parameters, labels or tuning
budget.

\textbf{H2, transferable utility:} a frozen consequence representation
supports new in-span goals with fewer grounding labels than fitting a
new scalar evaluator. Falsify the transfer claim when it depends on
evaluation-parent leakage, a task-specific encoder update or reward
properties not actually represented.

\textbf{H3, temporal composition:} trajectory modeling helps when useful
outcomes depend on multiple separated actions. Falsify the long-range
mechanism when gains remain unchanged under removal of dependencies, or
are explained by a larger context window alone.

\textbf{H4, discovered structure:} clustering or learned discrete codes
improve label efficiency, compositional transfer or interpretability
without harming rare outcomes. Falsify the benefit when continuous or
random-projection controls perform equally well, code usage collapses,
or clusters reflect surface task IDs rather than consequences.

\textbf{H5, deployable integration:} the native typed-slot route or a
sidecar improves actual task performance within the target cost and
preservation envelope. Falsify deployment readiness when the gain
vanishes after counting planning overhead, fails under the real
executor, or harms protected behavior beyond the preregistered margin.

\subsection{Minimal model matrix}\label{minimal-model-matrix}

\begin{longtable}[]{@{}
  >{\raggedright\arraybackslash}p{(\columnwidth - 4\tabcolsep) * \real{0.3333}}
  >{\raggedright\arraybackslash}p{(\columnwidth - 4\tabcolsep) * \real{0.3333}}
  >{\raggedright\arraybackslash}p{(\columnwidth - 4\tabcolsep) * \real{0.3333}}@{}}
\toprule\noalign{}
\begin{minipage}[b]{\linewidth}\raggedright
ID
\end{minipage} & \begin{minipage}[b]{\linewidth}\raggedright
Model / representation
\end{minipage} & \begin{minipage}[b]{\linewidth}\raggedright
Question isolated
\end{minipage} \\
\midrule\noalign{}
\endhead
\bottomrule\noalign{}
\endlastfoot
B0 & Legal random or simple heuristic choice & Environment difficulty
and trivial strategies \\
B1 & Frozen Rineya direct action/answer & Existing deployed baseline \\
B2 & Scalar trajectory Transformer & Strong non-factorized evaluator \\
B3 & Measured consequence features + goal readout & Benefit of explicit
factorization \\
B4 & Learned continuous consequences + readout & Added value of learned
geometry \\
B5 & B4 with train-only clustering or VQ & Added value of discrete
structure \\
B6 & Policy-conditioned successor features & Direct future accumulation
versus rollout \\
B7 & Native Rineya typed-slot adapter & Backbone-sharing benefit and
retention cost \\
O1 & Exact simulator evaluator on same candidates & Candidate-set
ceiling; not deployable baseline \\
\end{longtable}

Do not run all combinations immediately. B0, B2, B3 and O1 validate the
environment and contracts. B1 enters when native capture and action
execution are qualified. B4 is the central scientific arm. B5, B6 and B7
are conditional extensions. A scalar baseline must have a comparable
representation-learning budget; otherwise an apparent factorization gain
may come from additional pretraining.

\subsection{Staged gates}\label{staged-gates}

\textbf{G0: provenance and mechanics.} Pass schema, split, snapshot,
mask, finite-loss and deterministic replay checks. Pin source identity
and demonstrate capture agreement within a declared tolerance. Required
output: a qualification report with no native-quality claim.

\textbf{G1: environment and instrument.} The oracle solves the intended
class; a tiny diagnostic set can be fitted; a measured-feature model
learns nontrivial action dependence. Required output: failures
stratified by observation ambiguity, label validity and optimization. A
broken instrument blocks later architecture comparisons.

\textbf{G2: learned-representation utility.} Compare B2/B3/B4 on paired
held-out parents. Primary endpoint: actual utility or task success under
a fixed candidate set. Secondary endpoints: consequence fidelity, label
efficiency and runtime. Required output: effect estimate, uncertainty
interval and resource accounting.

\textbf{G3: reuse and discovery.} Run nested grounding budgets and
held-out goals. Add clustering only after the continuous result.
Required output: learning curves with identical labeled subsets and a
rare-failure audit. A pretty latent plot is insufficient.

\textbf{G4: composition and horizon.} Increase dependency distance,
horizon and DAG complexity separately. Required output: free-running
prediction and actual-control results, including model-exploitation and
support stress tests.

\textbf{G5: native adaptation and preservation.} Add typed
slots/adapters behind a mode switch. Required output: checkpoint
invariance where promised, acquired world-model skill, old-behavior
noninferiority, measured laptop memory and system latency.

\textbf{G6: optional policy improvement.} Distill a qualified planner or
add short imagined RL. Required output: real-environment improvement
beyond the non-RL controller at matched data and compute.
Simulation-only gains are not promotion evidence.

\subsection{Sample size and statistical
unit}\label{sample-size-and-statistical-unit}

The primary independent unit is the \textbf{parent problem or
environment family}, not every correlated branch, goal relabel or token.
Aggregate paired outcomes within parents, then use a parent-level paired
bootstrap or an appropriate clustered model. Preserve gained/lost cases
for binary results. Training-seed variability is a separate uncertainty
source; report seed-level results and a hierarchical summary where
justified. The statistical-precedence literature motivates avoiding
confident aggregate claims from fragile run summaries. {[}S34{]}

Use a pilot to estimate variance and choose a final sample size before
sealed evaluation. As an illustrative planning calculation, if paired
parent-level utility differences have standard deviation \(\sigma\), a
rough normal approximation for detecting a mean difference \(\delta\) at
5\% two-sided significance and 80\% power is

\[
n\approx\left(\frac{(1.96+0.84)\sigma}{\delta}\right)^2.
\]

This is a derived approximation, not a substitute for clustered, binary
or heavy-tailed simulation-based power analysis. With \(\sigma=0.25\)
and \(\delta=0.03\), it suggests roughly 545 independent parents. More
correlated branches do not create 545 independent parents.

Predeclare one primary comparison and a small set of secondary
questions. Correct or clearly label exploratory multiplicity. Freeze
model selection on development data. Sealed evaluation is not a
hyperparameter search tool.

\subsection{Practical effect, preservation and
cost}\label{practical-effect-preservation-and-cost}

A suggested contract requires a positive primary paired confidence
interval, a preregistered minimum meaningful effect, and no unacceptable
cost or retention violation. The actual effect threshold must be chosen
using the task scale and pilot variance. Configurations deliberately
leave final promotion thresholds unset until that decision; a null
threshold means \textbf{blocked}, not ``no threshold required.''

For preservation, select a noninferiority margin on protected
capabilities. Failure to reject a difference is not evidence of no harm.
Require the lower bound on new-minus-base performance to exceed the
allowed loss margin, while separately showing acquisition. A frozen base
weight hash is necessary for a strict frozen route but does not certify
behavior with the new controller. {[}I02{]}

Report peak allocated and reserved GPU memory, process-tree RAM, disk
bytes, model and adapter parameters, encoding time, candidate
generation, scoring, tool execution, total latency and failure retries.
Compare quality at equal wall-clock or token budgets as well as raw
accuracy. Never hide a larger candidate search behind ``same model
size.''

\subsection{Decision outcomes}\label{decision-outcomes}

Promote a mechanism when it clears the declared accuracy, uncertainty,
preservation and cost gates. Refine it when the failure is a diagnosed
contract or coverage issue. Pivot to the simpler baseline when the extra
mechanism provides no useful benefit at the agreed budget. Preserve a
dormant extension when prerequisites are unmet; do not describe it as
empirically disproved when it was never qualified.

\section{Resource engineering for an 8 GB
laptop}\label{resource-engineering-for-an-8-gb-laptop}

\subsection{Memory arithmetic before kernel
tuning}\label{memory-arithmetic-before-kernel-tuning}

For the historical 596,049,920-parameter subject, BF16 weights alone
require approximately 1.11 GiB. A rough 16-bytes-per-parameter
full-training accounting gives approximately 8.88 GiB before
activations, attention workspaces and runtime overhead. These are
derived storage estimates, not measured device usage. They explain why
full fine-tuning is not the first route in an 8 GB envelope.

A small sidecar can be trained from cached frozen representations with
no Rineya weights resident during that training phase. Native adapters
reduce optimizer state but still require activations for gradient
propagation. Quantization can reduce frozen-weight memory, but does not
eliminate context activations, and the exact quantized route must pass
capture and behavior parity checks. {[}S27, S28{]}

Caching strategy matters. A single BF16 vector of width 1,024 for each
of 100,000 events occupies about 195 MiB. Saving all 28 such layer
vectors would use about 5.34 GiB before metadata. Saving every token of
every context is much larger. These are illustrative dimensions; the
actual checkpoint configuration remains authoritative. Cache selected
action-boundary summaries first, and use chunked files rather than
loading the entire corpus into host memory.

\subsection{Three operating envelopes}\label{three-operating-envelopes}

\textbf{CPU mechanics:} no Rineya checkpoint, tiny synthetic states,
short horizons and small Transformer. This is the supplied reference
route. It validates accounting and interfaces, not language
intelligence.

\textbf{Frozen capture plus sidecar:} load the qualified Rineya route
only to extract bounded representations; flush cache shards; train the
small predictor separately. Begin at 512 tokens and one candidate
microbatch. Increase context or candidate count only after measuring
peak memory and timing.

\textbf{Native adapter experiment:} register trainable modules, enable
appropriate checkpointing, use a microbatch of one if necessary,
accumulate gradients and remeasure the exact precision/cache/mask
combination. Increase one dimension at a time. A previous short forward
pass does not authorize an 8,192-token backward pass.

Reserve explicit VRAM and host-memory headroom, and stop before the
operating system enters uncontrolled paging. The suggested native
configuration sets a conservative peak-allocation admission target
rather than assuming the whole advertised 8 GB is available.
Reserved-memory and process-tree-RAM checks remain necessary. Guard
against thermal throttling when comparing runs separated in time.

\subsection{Numerical qualification}\label{numerical-qualification}

Before capture, compare single-prefix and bulk causal predictions at the
same token positions. Check padding side, position IDs, attention masks,
cache use, dropout mode, precision and fused-kernel settings. Compare
both selected hidden vectors and raw next-token log probabilities where
accessible. A small score discrepancy can change candidate ranking near
a decision boundary.

Test save/resume continuity, immutable base hashes, optimizer-step
counts and accumulation scaling. Ensure that a skipped or failed batch
does not silently count as an update. Store maximum errors and the
declared tolerances, not only a pass/fail word. The historical Rineya
capture issue makes this gate especially important, but the report does
not assume that the historical issue remains unresolved today. {[}I01{]}

\section{Preservation and integration with existing Rineya
work}\label{preservation-and-integration-with-existing-rineya-work}

\subsection{Keep the base anchor and the scientific question
separate}\label{keep-the-base-anchor-and-the-scientific-question-separate}

The consequence model is a shared interface that can serve several
existing directions. It does not replace their controls. Hidden-state
research supplies candidate representations; TAPS/CDE supplies possible
compute-allocation actions; CSS and retention work supply adapter and
preservation mechanisms; post-training supplies trajectories and
optional policy updates. Integrate through contracts and shared logging
rather than immediately merging every experimental branch.

Adaptive recurrence is a later action family because it requires a
reliably executable recurrence mechanism and a measured benefit. A world
model cannot rescue a recurrence operator that does not produce useful
state changes. First compare the actual effects of additional recurrence
against ordinary extra tokens or a matched tool action. Then ask whether
the consequence predictor can choose when to use it.

For skill composition, keep acquisition, old-skill retention and joint
execution separate. A world-model adapter may retain old answers but
fail when two skills must interact in one plan. Evaluate held-out
compositions and real action sequences, not only separate skill test
scores. Low cosine similarity between parameter deltas does not certify
functional independence. {[}I02{]}

\subsection{Native mode isolation}\label{native-mode-isolation}

A \texttt{language} mode should preserve the original prompt path,
tokenizer, model weights and output head. A \texttt{consequence} mode
adds the registered projections/adapters and heads. Unit-test that
disabling the new mode restores the original forward behavior within the
required tolerance. If a shared adapter changes normal language
behavior, report it as a changed model, not an unchanged base because
only a small parameter fraction was trained.

Do not merge an adapter or delta into the anchor merely to simplify
packaging during research. Persist the anchor identity, adapter
identity, target modules, precision and merge recipe separately.
Quantized or low-precision deployment can alter small updates; qualify
the artifact that will actually run, not only a higher-precision
training copy.

\subsection{Teacher and data reuse}\label{teacher-and-data-reuse}

Existing post-training traces can be useful only when their action
boundaries, executor semantics and reward provenance are known.
Free-form reasoning text is not automatically an executable trajectory.
Tag unsupported retrospective labels as such. When teacher-generated
actions or labels are used, count teacher cost and retain the
possibility that multiple models share the same systematic error.

A distilled planner should be judged against real outcomes, not merely
its ability to reproduce a teacher's preferred plan. Retention replay
must come from allowed training data. Old sealed suites remain
evaluation-only, and current experiment failures must not be inserted
into training while preserving the same claimed test split.

\section{Engineering, safety and release
controls}\label{engineering-safety-and-release-controls}

\subsection{Repository organization and
interfaces}\label{repository-organization-and-interfaces}

Use an isolated experiment package with typed contracts, environment
adapters, representation models, utility heads, planners, data loaders,
evaluators and tests. Connect to the current repository through a narrow
\texttt{FrozenHistoryEncoder} protocol, rather than importing unverified
native internals throughout the experiment. A later integration pull
request should list which source files change, why they change and how
the ordinary language route remains covered.

The bundle includes standalone schemas, proposed configurations, a
machine-readable experiment registry, acceptance templates and an
implementation handoff. The reference package is intentionally smaller
than this target architecture. Its implementation-status table
identifies missing native, self-supervised and production capabilities.
Do not interpret a directory name or a passing synthetic test as
implementation of the entire research program.

\subsection{Data and checkpoint
reproducibility}\label{data-and-checkpoint-reproducibility}

A run manifest records code revision, dirty-tree state, checkpoint and
tokenizer hashes, environment version, split manifest, feature schema,
action registry, mask mode, seed, hardware/software versions, precision,
optimization parameters, effective batch, costs and artifact hashes.
Record both requested and actually executed budgets. Save atomic
checkpoint files and their dependent normalizers/readouts together.

Deterministic seeds are useful, but bitwise equivalence across different
devices or kernel versions should not be promised. Pin the tested
environment and quantify tolerances. Re-run an independent reproduction
before an important promotion. The supplied validation results describe
this container's CPU environment only.

\subsection{Safety boundaries for learned
utility}\label{safety-boundaries-for-learned-utility}

Tool permission and irreversible-action checks occur before learned
scoring. The model cannot authorize file deletion, payment, account
changes or external communication. Use mocked tools or a sandbox
environment. Constraint violations are separately reported and cannot be
averaged away by a high task reward.

Test malicious instructions embedded in retrieved documents, forged tool
outputs, stale evidence, missing prerequisites, repeated no-op loops and
candidates designed to exploit the scorer. Do not execute
model-generated code without a bounded sandbox. Tool outputs are data;
they must not change the source checkpoint, evaluation split or logging
policy.

\subsection{Release gates and honest
claims}\label{release-gates-and-honest-claims}

A research artifact may be released with a negative result. It may not
be labeled a trained Rineya world model without a pinned native
checkpoint evaluation. A model card should state training data, intended
tasks, failure modes, prohibited actions, calibration limits and whether
goals are in-span or extrapolative. A dataset card should describe
labels, provenance, privacy, split grouping and license review.

Third-party papers and code are linked through the resource register
rather than bundled indiscriminately. Verify the exact repository
license and revision before copying code or downloading datasets; a
paper's open-access status does not determine its software or dataset
license. The ZIP contains original planning material and original
reference code, not a redistribution of the cited implementations.

\section{Delivery roadmap and
ownership}\label{delivery-roadmap-and-ownership}

\subsection{A gated sequence, not a calendar
promise}\label{a-gated-sequence-not-a-calendar-promise}

The following effort blocks are planning estimates. One person can hold
several roles, but scientific decisions, data validation and
implementation should still receive separate reviews. Native hardware
availability, existing repository interfaces and the pilot's statistical
requirements determine actual duration.

\textbf{Block 1: approximately 3--5 focused work sessions.} Research
owner freezes the hypothesis, subject identity, environment, primary
metric and scope. Engineering owner qualifies schemas, split logic,
masks and snapshots. Reviewer checks that future information cannot
leak. Deliverable: G0 report and a tiny end-to-end synthetic run.

\textbf{Block 2: approximately 1--2 working weeks.} Implement the
measured-feature and scalar baselines, collect the first bounded parent
set, establish candidate-oracle performance and measure laptop capture.
Deliverable: G1 report, failure taxonomy, measured resource envelope and
finalized statistical contract. No self-supervised codebook is required
yet.

\textbf{Block 3: approximately 1--2 working weeks.} Add the learned
continuous outcome encoder/predictor and sparse goal readout. Run
B2/B3/B4 with matched data and budgets. Deliverable: G2 result and a
decision to promote, repair or prefer the simpler evaluator.

\textbf{Block 4: approximately 1--2 working weeks.} Test nested label
budgets, new goals and controlled clustering. Add direct successor
targets only if needed by the temporal failure analysis. Deliverable: G3
transfer and discovery report with negative controls.

\textbf{Block 5: approximately 1--2 working weeks.} Expand context and
dependency horizon, then introduce synchronous DAG tasks. Deliverable:
G4 report including causal visibility, joint outcomes, scheduling and
model-exploitation stress tests.

\textbf{Block 6: approximately 1--3 working weeks, conditional.} Port
the successful interface to native typed-slot adapters, run
acquisition/retention and hardware gates, then evaluate a distilled
policy or short imagined-RL extension. Deliverable: G5/G6 qualification
or an explicit reason to retain the sidecar. Visual modeling remains a
later independent campaign.

These blocks are not a commitment to continuous background work. The
priority is to complete a decisive comparison before expanding the
architecture. The roadmap and experiment registry include explicit
dependencies to prevent every idea from becoming simultaneous active
work.

\subsection{First implementation
handoff}\label{first-implementation-handoff}

The first agent or engineer should read \texttt{START\_HERE.md}, the
architecture and data contracts, then run the reference tests. Next,
inspect the current repository in read-only mode, identify the actual
checkpoint loader and capture interface, and prepare an integration diff
plan. Stop for clarification when the subject identity, sealed-data
policy, allowed paths or native budget is unresolved.

The first substantive pull request should contain contracts, a simulator
adapter, split checks, mask tests and an end-to-end measured-feature
baseline---not native training, recurrence, codebooks and imagined RL
all at once. A separate pull request should add frozen Rineya capture
after numerical qualification. Each should include its exact tests,
resource observations and limitations.

\subsection{Decision log at every
gate}\label{decision-log-at-every-gate}

Record what was expected, what was observed, which comparison isolates
the mechanism, what is uncertain and what next action follows. Maintain
rejected and deferred alternatives. Do not rewrite the original
hypothesis after seeing evaluation results without labeling the new
experiment. Useful stopping decisions are part of the research output.

\section{Risk register and diagnostic
responses}\label{risk-register-and-diagnostic-responses}

\begin{longtable}[]{@{}
  >{\raggedright\arraybackslash}p{(\columnwidth - 4\tabcolsep) * \real{0.3333}}
  >{\raggedright\arraybackslash}p{(\columnwidth - 4\tabcolsep) * \real{0.3333}}
  >{\raggedright\arraybackslash}p{(\columnwidth - 4\tabcolsep) * \real{0.3333}}@{}}
\toprule\noalign{}
\begin{minipage}[b]{\linewidth}\raggedright
Risk
\end{minipage} & \begin{minipage}[b]{\linewidth}\raggedright
Detect it
\end{minipage} & \begin{minipage}[b]{\linewidth}\raggedright
Response
\end{minipage} \\
\midrule\noalign{}
\endhead
\bottomrule\noalign{}
\endlastfoot
Future-information leakage & Causal perturbation and provenance tests &
Block training/evaluation; rebuild inputs \\
Representation collapse & Variance, rank and decoder tests & Fix
target/regularization before scaling \\
Hidden-state aliasing & Same observation, different history branches &
Improve belief/history representation \\
Proxy exploitation & Predicted gain rises, real gain falls with search &
Reduce search, add real anchors, revisit support \\
Goal entanglement & Preference-only goal changes physical prediction &
Separate task instruction from utility weights \\
Rare failure merged away & Cluster audit by failure type & Increase
resolution or retain explicit hazard features \\
Invalid action semantics & Same token, inconsistent executor behavior &
Version and parameterize actions \\
Policy-mixed successor target & Different continuations share one target
& Condition on policy/plan or label the mixture \\
Frozen-weight but worse system & End-to-end retention degrades &
Restrict routing or revert controller \\
Apparent efficiency from hidden cost & Full cost ledger changes ranking
& Recompare under total budget \\
Goal transfer outside feature span & New-property readout fails & Add
grounded features or collect targeted data \\
Native numerical mismatch & Prefix/bulk or resume parity failure & Stop
native extraction; repair mechanics \\
\end{longtable}

The highest-risk scientific shortcut is to call any useful latent reward
predictor a novel world model. Keep the claim tied to the tested outcome
variables and interventions. The highest-risk engineering shortcut is to
accept a perfect validation loss without proving that the evaluator saw
only deployment-available information.

\section{Future extensions and research
questions}\label{future-extensions-and-research-questions}

\subsection{Distributional consequences and reachable
plans}\label{distributional-consequences-and-reachable-plans}

After the deterministic linear-utility baseline succeeds, add stochastic
outcomes and nonlinear preferences. Compare direct successor
distributions with recursive latent rollouts and shared-scenario
sampling. Evaluate joint failure, tail risk and reachability. Do not
select an attractive predicted consequence unless a supported action or
policy can plausibly realize it. DiSPO supplies an especially relevant
conceptual reference for coupling future features to reachable behavior.
{[}S04{]}

\subsection{Learned action
vocabularies}\label{learned-action-vocabularies}

Manual operators provide a stable initial intervention vocabulary. A
later model can discover reusable macro-actions from trajectories, with
each code decoded into an executable, bounded option. Test whether
learned options transfer across tasks and retain meaningful
preconditions. Interpretability and compression are secondary to
reliable execution. An action code that cannot be executed consistently
is not a useful control primitive.

\subsection{Internal compute as an
environment}\label{internal-compute-as-an-environment}

Once native recurrence is qualified, model the consequences of
additional computation, checking, retrieval or stopping. Treat
hidden-state prediction as an internal computation model, and evaluate
whether it improves final correctness per actual latency. Compare
recurrence against ordinary extra generation and tool use. This is where
TAPS/CDE can become a downstream application rather than a duplicated
controller project.

\subsection{Visual version}\label{visual-version}

A later visual Rineya variant can attach a frozen vision encoder and
learn action-conditioned future visual features. Do not start by
generating pixels or reproducing internet-scale video pretraining.
V-JEPA 2, DINO-WM and V-JEPA 2.1 offer different
predictive-representation inspirations; the relevant experiment is
whether a compact action-conditioned model preserves task-critical
visual variables at the available budget. {[}S36, S37, S39{]}

Separate visual perception failures from dynamics errors by comparing
ground-truth state input, frozen visual features and learned visual
adaptation. Test unseen appearances with fixed dynamics and unseen
dynamics with similar appearances. A text-world success is not evidence
of physical-world competence.

\subsection{Cumulative skill
acquisition}\label{cumulative-skill-acquisition}

Investigate whether a consequence vocabulary supports adding a new goal
readout without changing old behavior, and whether adding genuinely new
transition knowledge can be localized in an adapter. Evaluate
composition of old and new operators in one plan. This links the project
to Rineya's broader aim of accumulating capabilities, but the world
model does not automatically solve catastrophic forgetting.

A particularly informative future study changes only the goal, then only
the dynamics, then only the observation rendering. These shifts test
three different forms of reuse. Calling all three ``generalization''
conceals which part of the architecture actually transfers.

\section{How to use the delivered
resources}\label{how-to-use-the-delivered-resources}

Read the executive decision and formal model first. Use the annotated
bibliography to understand the closest precedents and the limits of each
analogy. The \texttt{specs} directory is the implementation reference;
\texttt{configs}, \texttt{schemas} and \texttt{templates} turn the
report into explicit contracts. The \texttt{roadmap} directory records
dependencies and future gates. The \texttt{reference} package runs
without Rineya weights and demonstrates selected mechanics. Its own
README states what it does not implement.

Validation outputs are included separately from proposed experiment
results. The CPU demonstration is not a benchmark result and does not
justify a claim that learned consequence semantics, clustering,
long-context planning or skill retention have been solved. It is a
reproducible starting point for detecting mistakes before native
experiments become expensive.

The program's strongest claim, should later evidence support it, is
precise: \textbf{under a frozen data, compute and preservation budget, a
learned consequence representation improved held-out goal-conditioned
action selection or reduced the required grounding labels compared with
a matched scalar evaluator.} That is worth investigating. It is also
specific enough to fail, be refined and ultimately produce useful
knowledge.

\section{Glossary and notation}\label{glossary-and-notation}

\begin{longtable}[]{@{}
  >{\raggedright\arraybackslash}p{(\columnwidth - 2\tabcolsep) * \real{0.5000}}
  >{\raggedright\arraybackslash}p{(\columnwidth - 2\tabcolsep) * \real{0.5000}}@{}}
\toprule\noalign{}
\begin{minipage}[b]{\linewidth}\raggedright
Term
\end{minipage} & \begin{minipage}[b]{\linewidth}\raggedright
Meaning in this bundle
\end{minipage} \\
\midrule\noalign{}
\endhead
\bottomrule\noalign{}
\endlastfoot
\(H_t\), history & Observations and actions available at decision
time \\
\(b_t\), belief & Encoded history; not necessarily true environment
state \\
\(a_t\), action & Executable operator plus arguments, budget and
version \\
\(\phi_t\), consequence & Measured or learned description of a
transition \\
\(z_t\), latent & Learned consequence coordinates without assumed named
axes \\
\(g\), goal & Task/preference specification; roles recorded
separately \\
\(w_g\), utility weights & Readout for an explicitly represented feature
family \\
\(\psi^\pi\), successor features & Expected discounted consequences
under policy \(\pi\) \\
\(A\), candidate plan & Precommitted actions or an explicit conditional
policy tree \\
Counterfactual family & Branches, goals and variants from the same
underlying parent \\
Grounding & Independent task feedback connecting representations to
utility \\
In-span revaluation & Changing utility over already represented
consequences \\
Native qualification & Tests using the exact selected Rineya artifact
and runtime \\
Promotion & Meeting predeclared scientific, preservation and cost
gates \\
\end{longtable}

\clearpage

\section{Delivered reference mechanics and
validation}\label{delivered-reference-mechanics-and-validation}

\subsection{What was actually run}\label{what-was-actually-run}

The bundle's original CPU package was tested in the recorded
environment: Python 3.13.5, PyTorch 2.10.0+cpu and NumPy 2.3.5.
\textbf{60 tests passed.} They cover finite-horizon consequence
accounting, goal revaluation, mask conventions, future-action
invariance, parent grouping, train-only codebook fitting,
schema/temporal checks, simulator behavior, tensor shapes, finite
gradients and checkpoint roundtrip. Full test names and logs are
included. No GPU or native Rineya checkpoint was used.

A 40,325-parameter causal action-token Transformer was trained for 250
updates on an original KeyDoor toy dataset. The check-set mean-squared
feature error changed from 0.540924 to 0.002186. This is a mechanics
demonstration with repeated support in a tiny state space, not an
independent generalization benchmark. Its learned bottleneck is
supervised through measured features; it does not validate
self-supervised semantic discovery or clustering.

\subsection{An important observed
limitation}\label{an-important-observed-limitation}

The average error alone would be misleading. On the illustrated safe
sequence
\texttt{take\_key\ -\textgreater{}\ open\_door\ -\textgreater{}\ collect},
the model predicted a cumulative object-gained feature of approximately
\textbf{0.064}, although the exact simulator value was \textbf{1.000}.
Under the damage-averse goal, it scored this route at approximately
\textbf{-3.46}, although the exact value was \textbf{7.00}. The demo
therefore does \textbf{not} pass a strong per-plan fidelity gate merely
because its aggregate error fell.

This failure is deliberately preserved in the delivered evidence. It
illustrates why the proposed program requires prerequisite-stratified
checks, candidate regret and independent real outcomes, rather than only
an average reconstruction score. The code's mechanical tests passing
means the selected tested invariants held; it does not mean its learned
predictions are accurate enough for deployment. A later coverage or
training change should be a separately logged experiment, not a silent
rewrite of this run.

\subsection{Resource and artifact
checks}\label{resource-and-artifact-checks}

Three JSON schemas, four illustrative trajectory records, seven proposed
configurations and the 43-entry source register were checked by the
bundled validation script. Temporal and feature-dimension checks extend
the JSON Schema validation for the initial serial format. Production
provenance, native memory, stochastic calibration, parallel execution
and retention remain unmeasured.

The final archive includes a file manifest with SHA-256 hashes. The
manifest excludes itself to avoid a recursive hash definition. The
verification utility checks each listed file after extraction. These
integrity checks are reproducibility aids, not a security certification
or validation of scientific claims.

\subsection{Issues found during
construction}\label{issues-found-during-construction}

An initial unit-test run exposed a Python/PyTorch naming collision: a
submodule named \texttt{type} conflicted with the parent module's
dtype-conversion method. Renaming it to \texttt{type\_embedding} fixed
the forward path, after which the suite was rerun. This is recorded as a
resolved engineering issue, not hidden behind the final passing result.

The report and machine-readable status files distinguish implemented
reference mechanics from planned native work. The most important next
implementation task is still the matched scalar/feature baseline and
qualified Rineya capture, not an expansion into imagined RL.

\clearpage

\section{Source register}\label{source-register}

Sources accessed 5 October 2026. Full annotations and review-depth
labels are in \texttt{research/ANNOTATED\_READING.md}; machine-readable
metadata and BibTeX are included. Papers are evidence for their reported
settings, not measured Rineya outcomes.

\textbf{{[}S01{]}} Mastering diverse control tasks through world models.
Hafner, Danijar; Pasukonis, Jurgis; Ba, Jimmy; Lillicrap, Timothy. 2025;
Nature. \url{https://www.nature.com/articles/s41586-025-08744-2}

\textbf{{[}S02{]}} Successor Features for Transfer in Reinforcement
Learning. Barreto, Andre; others. 2017; NeurIPS.
\url{https://arxiv.org/abs/1606.05312}

\textbf{{[}S03{]}} Universal Successor Features Approximators. Borsa,
Diana; others. 2019; ICLR. \url{https://arxiv.org/abs/1812.07626}

\textbf{{[}S04{]}} Distributional Successor Features Enable Zero-Shot
Policy Optimization. Zhu, Chuning; Wang, Xinqi; Han, Tyler; Du, Simon
S.; Gupta, Abhishek. 2024; NeurIPS; arXiv v4 2025.
\url{https://arxiv.org/abs/2403.06328}

\textbf{{[}S05{]}} Successor Features for Efficient Multisubject
Controlled Text Generation. Cao, Meng; Fatemi, Mehdi; Cheung, Jackie Chi
Kit; Shabanian, Samira. 2023; arXiv preprint.
\url{https://arxiv.org/abs/2311.04921}

\textbf{{[}S06{]}} Decision Transformer: Reinforcement Learning via
Sequence Modeling. Chen, Lili; others. 2021; NeurIPS.
\url{https://arxiv.org/abs/2106.01345}

\textbf{{[}S07{]}} Offline Reinforcement Learning as One Big Sequence
Modeling Problem. Janner, Michael; Li, Qiyang; Levine, Sergey. 2021;
NeurIPS; author project. \url{https://trajectory-transformer.github.io/}

\textbf{{[}S08{]}} UniMASK: Unified Inference in Sequential Decision
Problems. Carroll, Micah; others. 2022; NeurIPS.
\url{https://arxiv.org/abs/2211.10869}

\textbf{{[}S09{]}} Masked Trajectory Models for Prediction,
Representation, and Control. Wu, Philipp; Majumdar, Arjun; Stone, Kevin;
Lin, Yixin; Mordatch, Igor; Abbeel, Pieter; Rajeswaran, Aravind. 2023;
ICML / PMLR. \url{https://proceedings.mlr.press/v202/wu23o.html}

\textbf{{[}S10{]}} Learning Action Representations for Reinforcement
Learning. Chandak, Yash; Theocharous, Georgios; Kostas, James; Jordan,
Scott; Thomas, Philip. 2019; ICML / PMLR.
\url{https://proceedings.mlr.press/v97/chandak19a.html}

\textbf{{[}S11{]}} Mastering Atari, Go, chess and shogi by planning with
a learned model. Schrittwieser, Julian; others. 2020; Nature; primary
arXiv version. \url{https://arxiv.org/abs/1911.08265}

\textbf{{[}S12{]}} DeepMDP: Learning Continuous Latent Space Models for
Representation Learning. Gelada, Carles; others. 2019; ICML / PMLR.
\url{https://proceedings.mlr.press/v97/gelada19a.html}

\textbf{{[}S13{]}} Learning Invariant Representations for Reinforcement
Learning without Reconstruction. Zhang, Amy; McAllister, Rowan;
Calandra, Roberto; Gal, Yarin; Levine, Sergey. 2021; ICLR.
\url{https://arxiv.org/abs/2006.10742}

\textbf{{[}S14{]}} Representation Learning with Contrastive Predictive
Coding. van den Oord, Aaron; Li, Yazhe; Vinyals, Oriol. 2018; arXiv
preprint. \url{https://arxiv.org/abs/1807.03748}

\textbf{{[}S15{]}} Neural Discrete Representation Learning. van den
Oord, Aaron; Vinyals, Oriol; Kavukcuoglu, Koray. 2017; NeurIPS.
\url{https://arxiv.org/abs/1711.00937}

\textbf{{[}S16{]}} VICReg: Variance-Invariance-Covariance Regularization
for Self-Supervised Learning. Bardes, Adrien; Ponce, Jean; LeCun, Yann.
2022; ICLR. \url{https://arxiv.org/abs/2105.04906}

\textbf{{[}S17{]}} Diversity is All You Need: Learning Skills without a
Reward Function. Eysenbach, Benjamin; others. 2019; ICLR.
\url{https://arxiv.org/abs/1802.06070}

\textbf{{[}S18{]}} Generalized Decision Transformer for Offline
Hindsight Information Matching. Furuta, Hiroki; Matsuo, Yutaka; Gu,
Shixiang Shane. 2021; arXiv preprint.
\url{https://arxiv.org/abs/2111.10364}

\textbf{{[}S19{]}} RUDDER: Return Decomposition for Delayed Rewards.
Arjona-Medina, Jose A.; others. 2019; NeurIPS.
\url{https://arxiv.org/abs/1806.07857}

\textbf{{[}S20{]}} Deep Reinforcement Learning from Human Preferences.
Christiano, Paul F.; others. 2017; NeurIPS.
\url{https://arxiv.org/abs/1706.03741}

\textbf{{[}S21{]}} Let's Verify Step by Step. Lightman, Hunter; others.
2023; arXiv preprint. \url{https://arxiv.org/abs/2305.20050}

\textbf{{[}S22{]}} When to Trust Your Model: Model-Based Policy
Optimization. Janner, Michael; Fu, Justin; Zhang, Marvin; Levine,
Sergey. 2019; NeurIPS. \url{https://arxiv.org/abs/1906.08253}

\textbf{{[}S23{]}} Conservative Q-Learning for Offline Reinforcement
Learning. Kumar, Aviral; others. 2020; NeurIPS.
\url{https://arxiv.org/abs/2006.04779}

\textbf{{[}S24{]}} A Distributional Perspective on Reinforcement
Learning. Bellemare, Marc G.; Dabney, Will; Munos, Remi. 2017; ICML.
\url{https://arxiv.org/abs/1707.06887}

\textbf{{[}S25{]}} Simple and Scalable Predictive Uncertainty Estimation
using Deep Ensembles. Lakshminarayanan, Balaji; Pritzel, Alexander;
Blundell, Charles. 2017; NeurIPS. \url{https://arxiv.org/abs/1612.01474}

\textbf{{[}S26{]}} Scaling Laws for Reward Model Overoptimization. Gao,
Leo; Schulman, John; Hilton, Jacob. 2022; arXiv preprint.
\url{https://arxiv.org/abs/2210.10760}

\textbf{{[}S27{]}} LoRA: Low-Rank Adaptation of Large Language Models.
Hu, Edward J.; others. 2022; ICLR.
\url{https://arxiv.org/abs/2106.09685}

\textbf{{[}S28{]}} QLoRA: Efficient Finetuning of Quantized LLMs.
Dettmers, Tim; others. 2023; NeurIPS.
\url{https://arxiv.org/abs/2305.14314}

\textbf{{[}S29{]}} RULER: What's the Real Context Size of Your
Long-Context Language Models?. Hsieh, Cheng-Ping; others. 2024; COLM.
\url{https://arxiv.org/abs/2404.06654}

\textbf{{[}S30{]}} LongBench: A Bilingual, Multitask Benchmark for Long
Context Understanding. Bai, Yushi; others. 2024; ACL.
\url{https://arxiv.org/abs/2308.14508}

\textbf{{[}S31{]}} TextWorld: A Learning Environment for Text-based
Games. Cote, Marc-Alexandre; others. 2018; arXiv preprint.
\url{https://arxiv.org/abs/1806.11532}

\textbf{{[}S32{]}} ALFWorld: Aligning Text and Embodied Environments for
Interactive Learning. Shridhar, Mohit; others. 2021; ICLR.
\url{https://arxiv.org/abs/2010.03768}

\textbf{{[}S33{]}} WebShop: Towards Scalable Real-World Web Interaction
with Grounded Language Agents. Yao, Shunyu; others. 2022; NeurIPS.
\url{https://arxiv.org/abs/2207.01206}

\textbf{{[}S34{]}} Deep Reinforcement Learning at the Edge of the
Statistical Precipice. Agarwal, Rishabh; others. 2021; NeurIPS.
\url{https://arxiv.org/abs/2108.13264}

\textbf{{[}S35{]}} Hindsight Experience Replay. Andrychowicz, Marcin;
others. 2017; NeurIPS. \url{https://arxiv.org/abs/1707.01495}

\textbf{{[}S36{]}} V-JEPA 2: Self-Supervised Video Models Enable
Understanding, Prediction and Planning. Assran, Mido; others. 2025;
arXiv preprint. \url{https://arxiv.org/abs/2506.09985}

\textbf{{[}S37{]}} DINO-WM: World Models on Pre-trained Visual Features
enable Zero-shot Planning. Zhou, Gaoyue; Pan, Hengkai; LeCun, Yann;
Pinto, Lerrel. 2024; arXiv preprint.
\url{https://arxiv.org/abs/2411.04983}

\textbf{{[}S38{]}} From Word to World: Can Large Language Models be
Implicit Text-based World Models?. Li, Yixia; others. 2025; arXiv
preprint; revision 2026. \url{https://arxiv.org/abs/2512.18832}

\textbf{{[}S39{]}} V-JEPA 2.1: Unlocking Dense Features in Video
Self-Supervised Learning. Mur-Labadia, Lorenzo; others. 2026; arXiv
preprint. \url{https://arxiv.org/abs/2603.14482}

\textbf{{[}S40{]}} Attention Is All You Need. Vaswani, Ashish; others.
2017; NeurIPS. \url{https://arxiv.org/abs/1706.03762}

\textbf{{[}S41{]}} torch.nn.Transformer documentation. \{PyTorch
contributors\}. 2026; official documentation; mutable.
\url{https://docs.pytorch.org/docs/2.10/generated/torch.nn.Transformer.html}

\textbf{{[}S42{]}} Qwen3 model documentation. \{Hugging Face
contributors\}. 2026; official documentation; mutable.
\url{https://huggingface.co/docs/transformers/model_doc/qwen3}

\textbf{{[}S43{]}} Gymnasium Env API. \{Farama Foundation
contributors\}. 2026; official documentation; mutable.
\url{https://gymnasium.farama.org/api/env/}

\subsection{Historical internal
references}\label{historical-internal-references}

\textbf{{[}I01{]}} Rineya Intelligence Research Handbook, 30 September
2026, study edition 1.0. Personal Library report; chapter 1, source
identity and laptop capture mechanics. The exact source checkpoint was
reported as 596,049,920 parameters; current identity and qualification
must be rechecked. Not a fresh repository audit.

\textbf{{[}I02{]}} Rineya: Architectures for Composable Skills, 21
September 2026. Personal Library report; sections 12--13 on paired
evaluation, preservation, compute reporting and sealed-data discipline.
These are historical research records, not external validation.

\textbf{{[}I03{]}} The present conversation, 5 October 2026. User
proposal: learn consequence/reward-related embeddings with action
tokens, attention, long-range and parallel structure, and
goal-conditioned evaluation. Concept origin; not empirical evidence.

\end{document}
